\documentclass[10pt,a4paper]{amsart}
\usepackage[margin=19mm]{geometry}
\usepackage{amsmath,amssymb,mathtools}
\usepackage[scr=rsfs,cal=euler]{mathalfa}
\usepackage{xfrac}
\usepackage[unicode,hidelinks,bookmarks=false]{hyperref}
\newcommand{\ie}{i.e.\ }
\newcommand{\cf}{cf.\ }
\newcommand{\resp}{respectively\ }
\newcommand{\Subsection}{\subsection}
\providecommand{\nobreakdash}{\nobreak\hbox{-}}
\setlength{\emergencystretch}{2em}
\multlinegap0pt
\usepackage{amssymb}
\usepackage{mathtools}
\usepackage[shortlabels]{enumitem}
\usepackage{xcolor}
\newtheorem{lemma}{Lemma}
\newcommand{\E}{\mathbb E}
\newcommand{\cA}{\mathcal A}
\newcommand{\ip}[2]{\left\langle #1,#2\right\rangle}
\newcommand{\one}{\mathbf{1}}
\newcommand{\taud}{\tau_d}
\title{Free-Probabilistic State Evolution and Random Matrix Discrepancy}
\author{August Y. Chen, Ahmed El Alaoui}
\date{Source: arXiv:2609.11836v1 (CC BY 4.0)}
\begin{document}
\maketitle

\noindent\emph{Source and licence.} Free-Probabilistic State Evolution and Random Matrix Discrepancy. Version arXiv:2609.11836v1 (CC BY 4.0), distributed by arXiv under the Creative Commons Attribution 4.0 International licence, http://creativecommons.org/licenses/by/4.0/, as recorded in the arXiv licence metadata for this exact version and shown on the abstract page https://arxiv.org/abs/2609.11836v1. Full licence notice: \texttt{licenses/arxiv-2609.11836-CC-BY-4.0.txt}.\par\smallskip

\noindent\textit{Editorial scope.} Counted unit: the article's lemma lem:adaptive-conditioning (source0.tex lines 1063--1133). The article's complete original proof of it is delivered with it (source0.tex lines 1135--1412, one \texttt{proof} environment of 278 lines). No mathematics is written, repaired or re-derived: the statement and the proof are the article's own lines, in the article's own notation, and no retained line is substituted or re-flowed.  Retained uncounted as the source-local context the proof needs: lines 1018--1020.  Nothing was omitted from any retained span, so the package carries no omission record.  No retained line contains a citation, so the wrapper carries no bibliography and no bibliography entry.  No retained line contains a figure environment, an image-inclusion command or drawing code, so no image is delivered and the package declares no asset dependency.  Editorial formatting only, in addition: an AMS \texttt{amsart} preamble that restates the article's own declarations and macros used by the retained lines, namely the environments \texttt{lemma} on separate counters, together with the standard packages those lines need.  The retained environments print in this document as Lemma 1 in order of appearance, whereas the article prints the same items with the numbering of its own full document.  The complete native source of the article as distributed in the arXiv e-print for this exact version is bundled unchanged as \texttt{sources/210054/source0.tex}.
\begin{equation}\label{eq:cov-goe}
  \E\,\taud(A_i M)\taud(A_i N)=\frac{2}{d^2}\ip{M}{N}_d\,.
\end{equation}

\begin{lemma}\label{lem:adaptive-conditioning}
Let
\[
  Q_T:\mathbb R^{T+1}\to H_n,
  \qquad
  Q_Te_t=q^t\,,~~~~~ t \le T 
\]
and
\[
  R_T:\mathbb R^{T+1}\to K_d,
  \qquad
  R_Te_t=r^t\,,~~~~~ t \le T
\]
where $(e_t)_{t=0}^{T}$ is the standard basis of $\mathbb R^{T+1}$.
Let
\[
  Y_T:=\mathcal A Q_T,
  \qquad
  X_T:=\mathcal A^*R_T\,.
\]
Assume that $Q_{T}$ and $R_T$ have full rank almost surely. Define
\[
  G_{Q}:=Q_T^*Q_T\,,
  \qquad
  G_{R}:=R_T^*R_T\,,
\]
and
\[
  P_{Q}:=Q_TG_{Q}^{-1}Q_T^*\,,
  \qquad
  P_{R}:=R_TG_{R}^{-1}R_T^*\,.
\] 
%
Let $\widetilde{\mathcal A}$ be a copy of $\mathcal{A}$ independent of
$\sigma(\mathcal{A}) \vee \mathscr{F}_{T+1}$. If
\[
  q=Q_T\beta+h,
  \qquad
  h\perp \operatorname{Ran}(Q_T)\,,
\]
where $q$ is $\mathscr F_{T+1}$-measurable, then, conditionally on
$\mathscr F_{T+1}$,
\begin{equation}\label{eq:adaptive-domain-vector}
  \mathcal A q
  \stackrel{\mathrm d}{=}
  Y_T\beta
  +
  R_{T}G_{R}^{-1}X_{T}^*h
  +
  P_{R}^\perp\widetilde{\mathcal A}h\,.
\end{equation}
Here and below, $R_{T-1}$ and $X_{T-1}$ denote $R_T$ and $X_T$
with their last columns deleted; when $T=0$, the corresponding terms are
absent.  If
\[
  k=R_{T-1}\gamma+\ell\,,
  \qquad
  \ell\perp \operatorname{Ran}(R_{T-1})\,,
\]
where $k$ is $\mathscr G_T$-measurable, then, conditionally on
$\mathscr G_T$,
\begin{equation}\label{eq:adaptive-range-vector}
  \mathcal A^* k
  \stackrel{\mathrm d}{=}
  X_{T-1}\gamma
  +
  Q_{T}G_{Q}^{-1}Y_{T}^*\ell
  +
  P_{Q}^\perp\widetilde{\mathcal A}^*\ell\,.
\end{equation}
\end{lemma}

\begin{proof}
For brevity write $Q=Q_T$, $R=R_T$, $Y=Y_T$, and $X=X_T$. We will first prove that conditionally on $\mathscr F_{T+1}$,
\begin{equation}\label{eq:adaptive-conditioning}
  \mathcal A
  \stackrel{\mathrm d}{=}
  Y G_Q^{-1}Q^* + R G_R^{-1}X^*P_Q^\perp + P_R^\perp\widetilde{\mathcal A}P_Q^\perp\,.
\end{equation}
This will prove \eqref{eq:adaptive-domain-vector}.  We prove the range-side formula~\eqref{eq:adaptive-range-vector} separately below by stopping the same conditioning argument at $\mathscr G_T$, before the final observation $x^T$ is revealed.

Let $\mathscr L(H_n,K_d)$ be the vector space of linear operators from $H_n$ to $K_d$ equipped with its Hilbert--Schmidt inner product: 
\[\langle B, C\rangle = \sum_{i=1}^n \ip{Bu_i}{Cu_i}_{d},\]
where $(u_i)$ is an orthonormal basis of $H_n$.
If $(F_j)$ is an orthonormal basis of $K_d$, then the coefficients
\[
  \langle \mathcal A u_i,F_j\rangle_d
\]
are independent centered Gaussian variables with common variance $2/d^2$. Indeed, we have $\mathcal A u_i \stackrel{\mathrm d}{=} A_i$, therefore each $\langle \cA u_i,F_j\rangle_d$ is a centered Gaussian. Moreover, we have by~\eqref{eq:cov-goe} that
\begin{align*}
\E\big[ \langle \cA u_i, F_j \rangle_d \langle \cA u_{i'}, F_{j'} \rangle_d \big] &= \frac1n \sum_{1 \le k, k' \le n} (u_i)_k (u_{i'})_{k'} \E\big[ \langle A_k, F_j \rangle_d \langle A_{k'}, F_{j'} \rangle_d \big] \\
&= \frac1n \one\{i=i', j=j' \} \sum_{k=1}^n \E\big[ \langle A_k, F_j \rangle_d^2 \big] = \frac2{d^2} \one\{i=i', j=j' \}\,,
\end{align*}
proving the assertion. Thus $\cA$ is a centered isotropic Gaussian element of $\mathscr L(H_n,K_d)$, up to the scalar variance $2/d^2$.

Next, we shall repeatedly use the following Gaussian fact. Let $E,F$ be finite-dimensional Hilbert spaces and let $B:E\to F$ be a centered isotropic Gaussian linear operator. If $h\in E$ is deterministic, then conditionally on $Bh$, the restriction of $B$ to $h^\perp$ is independent of $Bh$ and has an isotropic Gaussian law:
\begin{equation}\label{eq:one-step-domain}
  BP_{h^\perp}
  \stackrel{\mathrm d}{=}
  \widetilde B P_{h^\perp}\,,
\end{equation}
where $\widetilde B$ is an independent copy of $B$. Similarly, if $\ell\in F$ is deterministic, then conditionally on $B^*\ell$, 
\begin{equation}\label{eq:one-step-range}
  P_{\ell^\perp}B
  \stackrel{\mathrm d}{=}
  P_{\ell^\perp}\widetilde B\,.
\end{equation}
%
We now prove the adaptive conditioning statement. For $t=0,\ldots,T+1$, let
\[
  S_t:=\operatorname{span}\{q^0,\ldots,q^{t-1}\}\subset H_n\,,
  \qquad
  L_t:=\operatorname{span}\{r^0,\ldots,r^{t-1}\}\subset K_d\,,
\]
with the convention $S_0=L_0=\{0\}$. Let $P_{S_t}$ and $P_{L_t}$ denote the corresponding orthogonal projections.

We prove by induction that, conditionally on $\mathscr F_t$,
\begin{align}\label{eq:gaussianconditioning_induction_measurability}
  \mathcal A P_{S_t}
  \quad\text{and}\quad
  P_{L_t}\mathcal A
\end{align}
are $\mathscr F_t$-measurable, and the ``unrevealed block" satisfies
\begin{equation}\label{eq:residual-invariant}
  P_{L_t}^\perp \mathcal A P_{S_t}^\perp
  \stackrel{\mathrm d}{=}
  P_{L_t}^\perp \widetilde{\mathcal A} P_{S_t}^\perp\,,
\end{equation}
where $\widetilde{\mathcal A}$ is an independent copy of $\mathcal A$, independent of $\mathscr F_t$.

For $t=0$, this is immediate from the independence of $\mathscr F_0$ and $\mathcal A$, because $S_0=L_0=\{0\}$.

Now we assume~\eqref{eq:gaussianconditioning_induction_measurability},~\eqref{eq:residual-invariant} hold at time $t$, and decompose
\[
  q^t=P_{S_t}q^t+h^t\,, \qquad h^t:=P_{S_t}^\perp q^t\,.
\]
Therefore we may write
\begin{align*}
y^t=\mathcal A q^t = \mathcal A P_{S_t}q^t + P_{L_t}\mathcal A h^t + P_{L_t}^\perp \mathcal A h^t\,.
\end{align*}
Since $q^t$ and $S_t$ are $\mathscr F_t$-measurable, $h^t$ and $P_{S_t}q^t$ are $\mathscr F_t$-measurable. Thus by the inductive hypothesis, $\mathcal A P_{S_t}q^t$ and $P_{L_t}\mathcal A h^t$ are $\mathscr F_t$-measurable. Hence the only remaining random term conditionally on $\mathscr F_t$ is
\[
  P_{L_t}^\perp \mathcal A h^t
  =
  \big(P_{L_t}^\perp \mathcal A P_{S_t}^\perp\big)h^t\,.
\]

By the induction hypothesis, conditionally on $\mathscr F_t$, the operator $P_{L_t}^\perp \mathcal A P_{S_t}^\perp$ is a fresh isotropic Gaussian operator from $S_t^\perp$ to $L_t^\perp$. Applying the Gaussian conditioning fact~\eqref{eq:one-step-domain} to the vector $h^t$, we obtain that conditionally on $\mathscr F_t \vee \sigma( P_{L_t}^\perp \cA P_{S_t}^\perp h^t)= \mathscr F_t\vee\sigma(y^t) = \mathscr G_t$,
\begin{align}\label{eq:gaussianconditioning_induction_intermediatefreshgaussian}
  P_{L_t}^\perp \mathcal A P_{S_{t+1}}^\perp
  \stackrel{\mathrm d}{=}
  P_{L_t}^\perp \widetilde{\mathcal A}P_{S_{t+1}}^\perp\,,
\end{align}
with $\widetilde{\mathcal A}$ independent of $\mathscr G_t$. (The equality of the $\sigma$-algebras follows by our earlier remarks.) Moreover, $\mathcal A P_{S_{t+1}}$ is now $\mathscr G_t$-measurable, because $\mathcal A P_{S_t}$ is  $\mathscr F_t$-measurable and $y^t=\mathcal A q^t$ reveals the action of $\mathcal A$ in the new direction $q^t$. The operator $P_{L_t}\mathcal A$ is $\mathscr F_t$- and therefore $\mathscr G_t$-measurable.

Next, decompose
\[
  r^t=P_{L_t}r^t+\ell^t,
  \qquad
  \ell^t:=P_{L_t}^\perp r^t\,.
\]
Therefore
\begin{align*}
x^t = \mathcal A^*r^t = \mathcal A^*P_{L_t}r^t + P_{S_{t+1}} \cA^* \ell^t  + P_{S_{t+1}}^\perp \cA^* \ell^t\,.
\end{align*}
Recall that $r^t$ and $P_{L_t}\mathcal A$ are $\mathscr G_t$-measurable. Thus $\mathcal A^*P_{L_t}r^t$ is $\mathscr G_t$-measurable as well. In addition, the projection of $\mathcal A^*\ell^t$ onto $S_{t+1}$ is $\mathscr G_t$-measurable: for every $s\in S_{t+1}$,
\[
  \langle s,\mathcal A^*\ell^t\rangle_{n}
  =
  \langle \mathcal A s,\ell^t\rangle_{d}\,,
\]
and $\mathcal A s$, as well as $r^t, L_t$ and therefore $\ell^t$, are all $\mathscr G_t$-measurable. Therefore the only remaining random term conditionally on $\mathscr G_t$ is
\[
  P_{S_{t+1}}^\perp\mathcal A^*\ell^t
  =
  \big(P_{L_t}^\perp\mathcal A P_{S_{t+1}}^\perp\big)^*\ell^t\,.
\]

As shown by~\eqref{eq:gaussianconditioning_induction_intermediatefreshgaussian}, conditioned on $\mathscr G_t$, $P_{L_t}^\perp \widetilde{\mathcal A} P_{S_{t+1}}^\perp$ is a fresh isotropic Gaussian operator from $S_{t+1}^\perp$ to $L_t^\perp$. Applying the Gaussian conditioning fact~\eqref{eq:one-step-range} in the range direction $\ell^t$, we have that conditionally on
$\mathscr G_t\vee \sigma(\cA^* \ell^t) = \mathscr G_t\vee\sigma(x^t) = \mathscr F_{t+1}$,
\[
  P_{L_{t+1}}^\perp \mathcal A P_{S_{t+1}}^\perp
  \stackrel{\mathrm d}{=}
  P_{L_{t+1}}^\perp \widetilde{\mathcal A}P_{S_{t+1}}^\perp\,,
\]
with $\widetilde{\mathcal A}$ independent of $\mathscr F_{t+1}$.

Furthermore, $P_{L_{t+1}}\mathcal A$ is now $\mathscr F_{t+1}$-measurable. Indeed, $P_{L_t}\mathcal A$ is $\mathscr G_t$ and therefore $\mathscr F_{t+1}$-measurable, hence $\mathcal A^*P_{L_t}r^t$ is $\mathscr F_{t+1}$-measurable. Since $x^t=\mathcal A^*r^t$ together with $\mathcal A^*P_{L_t}r^t$ determines $\mathcal A^*\ell^t$, it follows that the component of $\mathcal A$ in the new range direction $\ell^t$, and hence $P_{L_{t+1}}\mathcal A$, is $\mathscr F_{t+1}$-measurable. As shown earlier, the operator $\mathcal A P_{S_{t+1}}$ is $\mathscr G_t$- and therefore $\mathscr F_{t+1}$-measurable. This completes the induction, proving~\eqref{eq:gaussianconditioning_induction_measurability},~\eqref{eq:residual-invariant}.

We now take $t=T+1$ in~\eqref{eq:residual-invariant}. Note we have
\[
  S_{T+1}=\operatorname{Ran}(Q)\,,
  \qquad
  L_{T+1}=\operatorname{Ran}(R)\,.
\]
Thus, conditionally on $\mathscr F_{T+1}$, we have
\begin{align}\label{eq:residual-invariant-result}
  P_R^\perp \mathcal A P_Q^\perp
  \stackrel{\mathrm d}{=}
  P_R^\perp\widetilde{\mathcal A}P_Q^\perp\,,
\end{align}
with $\widetilde{\mathcal A}$ independent of $\mathscr F_{T+1}$.

We now prove~\eqref{eq:adaptive-conditioning},~\eqref{eq:adaptive-domain-vector}. We first identify the deterministic, already revealed part of $\cA$. Since $Y=\mathcal A Q$,
\[
  \mathcal A P_Q
  =
  \mathcal A QG_Q^{-1}Q^*
  =
  YG_Q^{-1}Q^*\,.
\]
Similarly, since $X=\mathcal A^*R$, we have $X^*=R^*\mathcal A$, and hence
\[
  P_R\mathcal A
  =
  RG_R^{-1}R^*\mathcal A
  =
  RG_R^{-1}X^*\,.
\]
Therefore
\[
  P_R\mathcal A P_Q^\perp
  =
  RG_R^{-1}X^*P_Q^\perp\,.
\]
Using the orthogonal block decomposition
\[
  \mathcal A
  =
  \mathcal A P_Q
  +
  P_R\mathcal A P_Q^\perp
  +
  P_R^\perp\mathcal A P_Q^\perp\,,
\]
we obtain from~\eqref{eq:residual-invariant-result} that conditionally on $\mathscr F_{T+1}$,
\[
  \mathcal A
  \stackrel{\mathrm d}{=}
  YG_Q^{-1}Q^*
  +
  RG_R^{-1}X^*P_Q^\perp
  +
  P_R^\perp\widetilde{\mathcal A}P_Q^\perp\,.
\]
This proves \eqref{eq:adaptive-conditioning}.

Now let
\[
  q=Q\beta+h\,,
  \qquad
  h\perp\operatorname{Ran}(Q)\,.
\]
Then $Q^*h=0$ and $P_Q^\perp q=h$. Applying \eqref{eq:adaptive-conditioning} to $q$, we obtain
\[
\begin{aligned}
  \mathcal A q
  &\stackrel{\mathrm d}{=}
  YG_Q^{-1}Q^*(Q\beta+h)
  +
  RG_R^{-1}X^*P_Q^\perp(Q\beta+h)
  +
  P_R^\perp\widetilde{\mathcal A}P_Q^\perp(Q\beta+h) \\
  &=
  Y\beta
  +
  RG_R^{-1}X^*h
  +
  P_R^\perp\widetilde{\mathcal A}h\,.
\end{aligned}
\]
This proves \eqref{eq:adaptive-domain-vector}.

Finally, we prove~\eqref{eq:adaptive-range-vector}. Let
\[
 R_-:=R_{T-1}\,,\qquad X_-:=X_{T-1}\,,
 \qquad G_{R,-}:=R_-^*R_-\,,
 \qquad P_{R,-}:=R_-G_{R,-}^{-1}R_-^*\,.
\]
When $T=0$, all four objects are interpreted as the corresponding zero maps, and the terms containing them are absent. We stop the alternating conditioning induction after revealing $y^T=\cA q^T$, and before revealing $x^T=\cA^*r^T$.  At this point the known domain and range spaces are $\operatorname{Ran}(Q_T)$ and $\operatorname{Ran}(R_-)$, respectively. The same calculation used for \eqref{eq:adaptive-conditioning} gives, conditionally on $\mathscr G_T$,
\begin{equation}\label{eq:adaptive-conditioning-range-stop}
 \cA
 \stackrel{\mathrm d}{=}
 YG_Q^{-1}Q^*
 +R_-G_{R,-}^{-1}X_-^*P_Q^\perp
 +P_{R,-}^\perp\widetilde{\cA}P_Q^\perp\,,
\end{equation}
where $\widetilde{\cA}$ is independent of $\mathscr G_T$.  In our current notation, we may write
\[
  k=R_-\gamma+\ell\,, \qquad \ell\perp\operatorname{Ran}(R_-)\,.
\]
Taking adjoints in~\eqref{eq:adaptive-conditioning-range-stop} yields
\[
  \mathcal A^*
  \stackrel{\mathrm d}{=}
  QG_Q^{-1}Y^*
  +P_Q^\perp X_-G_{R,-}^{-1}R_-^*
  +P_Q^\perp\widetilde{\mathcal A}^*P_{R,-}^\perp\,.
\]
Therefore
\[
\begin{aligned}
  \mathcal A^*k
  &\stackrel{\mathrm d}{=}
  QG_Q^{-1}Y^*(R_-\gamma+\ell)
  +
  P_Q^\perp X_-G_{R,-}^{-1}R_-^*(R_-\gamma+\ell)
  +
  P_Q^\perp\widetilde{\mathcal A}^*P_{R,-}^\perp(R_-\gamma+\ell)  \\
  &=
  QG_Q^{-1}Y^*R_-\gamma
  +
  QG_Q^{-1}Y^*\ell
  +
  P_Q^\perp X_-\gamma
  +
  P_Q^\perp\widetilde{\mathcal A}^*\ell\,.
\end{aligned}
\]
Now observe that
\[
  Y^*R_-=Q^*\mathcal A^*R_-=Q^*X_-\,.
\]
This implies
\[
  QG_Q^{-1}Y^*R_-\gamma
  =
  QG_Q^{-1}Q^*X_-\gamma
  =
  P_QX_-\gamma\,.
\]
Combining with the earlier display thus gives
\begin{align}
\mathcal A^*k
  &\stackrel{\mathrm d}{=}
  P_QX_-\gamma
  +
  P_Q^\perp X_-\gamma
  +
  QG_Q^{-1}Y^*\ell
  +
  P_Q^\perp\widetilde{\mathcal A}^*\ell \notag \\
  &= X_-\gamma
  +
  QG_Q^{-1}Y^*\ell
  +
  P_Q^\perp\widetilde{\mathcal A}^*\ell\,. \notag
\end{align}
This proves \eqref{eq:adaptive-range-vector}, completing the proof of Lemma \ref{lem:adaptive-conditioning}.
\end{proof}

\end{document}
